\chapter{Customer Needs}

\section{Session Overview}

Customer needs identification is the bridge between product planning and technical specification. In this session, your group will gather customer-facing evidence, translate raw statements into needs, and organise those needs so they can be used in Session 8 for product specifications.

By the end of this session, you should be able to:
\begin{enumerate}
    \item generate effective survey or interview prompts;
    \item translate customer statements into clear needs statements;
    \item explain the value of lead users and extreme users; and
    \item compile a structured customer-needs table for your hydroponic product.
\end{enumerate}

\section{Key Terms}

The following terms are used throughout the chapter:
\begin{itemize}
    \item \textbf{Customer statement}: a direct quote or paraphrase describing a user's experience, frustration, or suggestion.
    \item \textbf{Needs statement}: a concise description of what the product must enable, reduce, or improve.
    \item \textbf{Lead user}: a user who experiences needs earlier than the mainstream market and may reveal emerging requirements.
    \item \textbf{Extreme user}: a user operating at the edge of normal use conditions and therefore likely to expose hidden limitations.
    \item \textbf{Variable type}: a broad category used to classify a need, such as functional, usability, cost, safety, or maintainability.
\end{itemize}

\section{Why Customer Needs Matter}

Customer needs help your group avoid designing only from assumptions. They provide a structured way to connect the mission statement from Sessions 5 and 6 to measurable engineering work. For hydroponic products, needs may involve fluid handling, maintenance, structural convenience, sensing accuracy, safety, or cost. The stronger your needs list, the easier it becomes to define meaningful specifications later.

\section{Pre-Class Activity 1: Generate Survey or Interview Prompts}

Before class, draft at least three questions or prompts that could help your group collect useful user feedback.

\subsection{Survey Design Guidelines}

Good questions should be:
\begin{itemize}
    \item clear and easy to understand;
    \item neutral rather than leading;
    \item connected to a real usage situation; and
    \item useful for revealing frustrations, preferences, or desired improvements.
\end{itemize}

Table~\ref{tab:session7-question-template} can be used to record your prompts.

\begin{table}[htbp]
\centering
\caption{Template for survey or interview prompts}
\label{tab:session7-question-template}
\renewcommand{\arraystretch}{1.2}
\begin{tabularx}{\textwidth}{|>{\raggedright\arraybackslash}p{0.08\textwidth}|>{\raggedright\arraybackslash}X|>{\raggedright\arraybackslash}p{0.28\textwidth}|}
\hline
\textbf{No.} & \textbf{Question or prompt} & \textbf{Why it is useful} \\
\hline
1 & How often do you check nutrient level or water level in your hydroponic system? & Reveals monitoring frequency and maintenance burden. \\
\hline
2 & What part of maintaining your hydroponic system is most frustrating? & Helps identify pain points and improvement opportunities. \\
\hline
3 & What would make the system easier to clean, refill, or monitor? & Encourages users to describe desired features directly. \\
\hline
\end{tabularx}
\end{table}

\section{Pre-Class Activity 2: Translate Customer Statements into Needs}

After collecting observations, imagined use cases, or preliminary feedback, convert each customer statement into a needs statement.

\subsection{Translation Guidance}

A strong needs statement should:
\begin{itemize}
    \item focus on what the product should achieve, not on the exact solution;
    \item remain concise and specific;
    \item avoid embedded assumptions about technology; and
    \item reflect the customer's underlying problem.
\end{itemize}

Table~\ref{tab:session7-needs-example} shows a hydroponic example.

\begin{table}[htbp]
\centering
\caption{Example of translating hydroponic customer statements into needs}
\label{tab:session7-needs-example}
\renewcommand{\arraystretch}{1.2}
\begin{tabularx}{\textwidth}{|>{\raggedright\arraybackslash}X|>{\raggedright\arraybackslash}X|>{\raggedright\arraybackslash}p{0.2\textwidth}|}
\hline
\textbf{Customer statement} & \textbf{Needs statement} & \textbf{Variable type} \\
\hline
``I never know when the nutrient tank is close to empty.'' & The system should provide a clear indication of nutrient level. & Functional \\
\hline
``Cleaning the pipes takes too long.'' & The system should be easy to clean and maintain. & Maintainability \\
\hline
``I worry that leaks will damage the floor near my plants.'' & The system should minimise leakage during normal operation. & Safety \\
\hline
``I want something simple because I am new to hydroponics.'' & The system should be easy for a beginner to understand and operate. & Usability \\
\hline
\end{tabularx}
\end{table}

\section{Pre-Class Activity 3: Consider Lead Users and Extreme Users}

Lead users and extreme users can reveal needs that ordinary users may not express clearly.

\begin{itemize}
    \item A \textbf{lead user} might be an experienced hydroponic grower who has already created workarounds for monitoring or dosing problems.
    \item An \textbf{extreme user} might be someone operating in a very small apartment, in a very hot environment, or under strict cleaning constraints.
\end{itemize}

Write a short note explaining whether these user groups are important for your product and why. If they are relevant, identify what special insight they may provide.

\section{Optional Reflection: Needs and Innovation}

Use this prompt for a short discussion or reflective paragraph: How can a carefully identified unmet need lead to a novel product concept? Which hidden or latent need in your hydroponic project might create the biggest design opportunity?

\section{In-Class Activity: Build the Customer Needs Table}

In class, merge your individual prompts and needs statements into one group table.

\subsection{Required Deliverable}

Your group should prepare a table containing:
\begin{itemize}
    \item the agreed survey or interview prompts;
    \item customer statements grouped by themes such as typical use, likes, dislikes, and suggested improvements;
    \item the interpreted needs statements; and
    \item the variable type for each interpreted need.
\end{itemize}

Table~\ref{tab:session7-group-template} provides the recommended format.

\begin{table}[htbp]
\centering
\caption{Compiled customer statements and interpreted needs for the selected hydroponic product}
\label{tab:session7-group-template}
\renewcommand{\arraystretch}{1.2}
\begin{tabularx}{\textwidth}{|>{\raggedright\arraybackslash}p{0.16\textwidth}|>{\raggedright\arraybackslash}X|>{\raggedright\arraybackslash}X|>{\raggedright\arraybackslash}p{0.16\textwidth}|}
\hline
\textbf{Prompt or theme} & \textbf{Customer statement} & \textbf{Interpreted need} & \textbf{Variable type} \\
\hline
Typical use &  &  &  \\
\hline
Like &  &  &  \\
\hline
Dislike &  &  &  \\
\hline
Suggested improvement &  &  &  \\
\hline
\end{tabularx}
\end{table}

If your group uses an LLM to help organise responses, document the prompts used and verify that the grouped output still reflects the actual evidence collected.

\section{Summary and Next Steps}

At the end of this session, your group should have a clear customer-needs table and a short note about whether lead or extreme users matter for the project. In Session 8, these needs will be translated into measurable product specifications.
